% !TEX program = xelatex
% Шаблон научной статьи НИТУ «МИСиС»
% Заполните блоки ниже и скомпилируйте: xelatex article.tex

\documentclass{misisarticle}

% --- Классификационные коды (оставьте пустым {}, если не требуются) ---
\grnti{}
\asjc{}
\oecd{}

% --- Название ---
\title{Формирование доказательного контекста для больших языковых моделей на основе графа знаний}

% --- Автор ---
\misisauthor{А.~В.~Берендеев}{кафедра инженерной кибернетики НИТУ «МИСиС»}{Москва, Россия}{m2303742@edu.misis.ru}

% --- Аннотация и ключевые слова ---
\abstract{%
  Рассматривается задача формирования проверяемого доказательного контекста (evidence) для больших языковых моделей при ответах на запросы и проверке утверждений с опорой на граф знаний. Источником знаний выступает Wikidata, загруженная в Neo4j в виде триплетов «субъект—предикат—объект». Предлагается метод построения evidence-контекста, включающий извлечение релевантного подграфа по сущностям запроса, фильтрацию шумовых и противоречивых фактов и преобразование триплетов в краткие предложения на естественном языке. Полученный контекст используется в RAG/GraphRAG-пайплайне как основание для генерации ответа или верификации утверждения. Подход ориентирован на полностью локальное выполнение и может применяться в корпоративных системах, где важны воспроизводимость и контроль источников знаний.
}

\keywords{%
  GraphRAG, большие языковые модели, верификация фактов.
}

\begin{document}

\maketitle

% ============================================================
\section{Введение}

Первый абзац раздела начинается с красной строки (отступ 0{,}51~см). Текст выравнивается по ширине, шрифт Times New Roman, 10~pt.

Второй абзац. Ссылки на литературу оформляйте так~\cite{example2024}. Для нескольких источников:~\cite{example2024, book2023}.

% ============================================================
\section{Основная часть}

\subsection{Подраздел}

Подзаголовки второго уровня набираются \textit{курсивом}.

Пример маркированного списка:
\begin{itemize}
  \item первый пункт;
  \item второй пункт;
  \item третий пункт.
\end{itemize}

Пример формулы в строке: мера Жаккара $IoU = \dfrac{|A \cap B|}{|A \cup B|}$.

Нумерованная формула (ссылка в тексте: формула~\eqref{eq:f1}):
\begin{equation}
  F_1 = 2 \cdot \frac{P \cdot R}{P + R}
       = \frac{2TP}{2TP + FP + FN},
  \label{eq:f1}
\end{equation}
где $P$~--- precision, $R$~--- recall, $TP$, $FP$, $FN$~--- число истинно и ложно положительных/отрицательных детекций.

% --- Рисунок ---
\begin{figure}[htbp]
  \centering
  % Замените заглушку на свой файл:
  % \includegraphics[width=0.75\linewidth]{figures/figure}
  \fbox{\parbox[c][3cm][c]{0.9\columnwidth}{\centering Placeholder для рисунка}}
  \caption{Пример иллюстрации. Замените заглушку на свой файл из папки figures.}
  \label{fig:example}
\end{figure}

% --- Таблица ---
\begin{table}[htbp]
  \caption{Пример таблицы с результатами}
  \label{tab:example}
  \centering
  \begin{tabular}{lcc}
    \toprule
    \textbf{Метод} & \textbf{Precision} & \textbf{Recall} \\
    \toprule
    Метод~A & 0.85 & 0.78 \\
    Метод~B & 0.91 & 0.83 \\
    \bottomrule
  \end{tabular}
\end{table}

% ============================================================
\section{Заключение}

Кратко сформулируйте основные выводы работы.

% ============================================================
% Список литературы — вариант 1: вручную
\begin{thebibliography}{9}

\bibitem{example2024} Author A., Author B. Title of the paper. \textit{Journal Name}, vol.~1, no.~2, pp.~100--110, 2024.

\bibitem{book2023} Author C. \textit{Book Title}. Publisher, City, 2023.

\end{thebibliography}

\end{document}
